% ============================================================================
% Simplified Chinese UI and structural strings for the One Math Book series.
% Loaded by styles/onemath.sty when \booklang is zh. Build with XeLaTeX
% (see the *_zh.tex entries). The places where Chinese puts the number inside
% the noun (第3章, 第一部分, the solution header) are in the zh block at the
% end of styles/onemath.sty.
% ============================================================================

% Theorem / statement environment titles
\newcommand{\omnameDefinition}{定义}
\newcommand{\omnameTheorem}{定理}
\newcommand{\omnameProposition}{命题}
\newcommand{\omnameLemma}{引理}
\newcommand{\omnameCorollary}{推论}
\newcommand{\omnameMethod}{方法}
\newcommand{\omnameExample}{例}
\newcommand{\omnameNotation}{记号}
\newcommand{\omnameRemark}{注}
\newcommand{\omnameExercise}{练习}
\newcommand{\omnameExercises}{练习}
\newcommand{\omnameProblem}{问题}
\newcommand{\omnameProblems}{问题}
\newcommand{\omnameProof}{证明}

% Admitted results and solutions appendix headers
\newcommand{\omadmittedtext}{在本阶段不加证明地承认。}
\newcommand{\omsolutionof}{的解答}
\newcommand{\ompage}{页}
\newcommand{\omnameGotoSolution}{解答见}

% Misc text macros
\newcommand{\st}{\text{，使得 }}

% Cover / title page
\newcommand{\bookmaintitle}{One Math Book}
\newcommand{\booktagline}{一本书，学完所有数学}
\newcommand{\bookdescription}{从幼儿园到本科毕业的数学}
\newcommand{\bookcontributors}{One Math Book 贡献者}
\newcommand{\bookversionlabel}{版本}

% Colophon
\newcommand{\omcolophonsource}{源代码、问题反馈与贡献：}
\newcommand{\omcolophonlicense}{除另有说明外，本书采用 Creative Commons CC BY-NC-SA 4.0 许可协议：}
\newcommand{\omcolophoncredit}{署名请注明：One Math Book, One Course (one-course.com)。}

% Solutions appendix part title
\newcommand{\omsolutionspart}{练习解答}

% Part titles (Grades 1–12; parallel to English Grade N)
\@namedef{omstr@part.grade1}{一年级（6--7 岁）}
\@namedef{omstr@part.grade2}{二年级（7--8 岁）}
\@namedef{omstr@part.grade3}{三年级（8--9 岁）}
\@namedef{omstr@part.grade4}{四年级（9--10 岁）}
\@namedef{omstr@part.grade5}{五年级（10--11 岁）}
\@namedef{omstr@part.grade6}{六年级（11--12 岁）}
\@namedef{omstr@part.grade7}{七年级（12--13 岁）}
\@namedef{omstr@part.grade8}{八年级（13--14 岁）}
\@namedef{omstr@part.grade9}{九年级（14--15 岁）}
\@namedef{omstr@part.grade10}{十年级（15--16 岁）}
\@namedef{omstr@part.grade11}{十一年级（16--17 岁）}
\@namedef{omstr@part.grade12}{十二年级（17--18 岁）}

% Part titles (Bachelor Years 1–3)
\@namedef{omstr@part.bachelor1}{本科一年级（18--19 岁）}
\@namedef{omstr@part.bachelor2}{本科二年级（19--20 岁）}
\@namedef{omstr@part.bachelor3}{本科三年级（20--21 岁）}

% Plural names + \omnameChapter, used by the \crefname declarations in
% onemath.sty (Chinese has no plural: the same word).
\newcommand{\omnameDefinitions}{定义}
\newcommand{\omnameTheorems}{定理}
\newcommand{\omnamePropositions}{命题}
\newcommand{\omnameLemmas}{引理}
\newcommand{\omnameCorollarys}{推论}
\newcommand{\omnameMethods}{方法}
\newcommand{\omnameExamples}{例}
\newcommand{\omnameNotations}{记号}
\newcommand{\omnameRemarks}{注}
\newcommand{\omnameChapter}{第}
\newcommand{\omnameChapters}{第}

% LaTeX's own structural names (no babel for Chinese).
\renewcommand{\chaptername}{第}
\renewcommand{\partname}{第}
\renewcommand{\contentsname}{目录}
\renewcommand{\appendixname}{附录}
\renewcommand{\indexname}{索引}
% No babel module is loaded for Chinese, so \today would print "November 17, 2026"
\renewcommand{\today}{\the\year\,年\,\the\month\,月\,\the\day\,日}
